%%=============================================================================
%% Methodologie
%%=============================================================================

\chapter{\IfLanguageName{dutch}{Methodologie}{Methodology}}%
\label{ch:methodologie}

%% TODO: In dit hoofstuk geef je een korte toelichting over hoe je te werk bent
%% gegaan. Verdeel je onderzoek in grote fasen, en licht in elke fase toe wat
%% de doelstelling was, welke deliverables daar uit gekomen zijn, en welke
%% onderzoeksmethoden je daarbij toegepast hebt. Verantwoord waarom je
%% op deze manier te werk gegaan bent.
%% 
%% Voorbeelden van zulke fasen zijn: literatuurstudie, opstellen van een
%% requirements-analyse, opstellen long-list (bij vergelijkende studie),
%% selectie van geschikte tools (bij vergelijkende studie, "short-list"),
%% opzetten testopstelling/PoC, uitvoeren testen en verzamelen
%% van resultaten, analyse van resultaten, ...
%%
%% !!!!! LET OP !!!!!
%%
%% Het is uitdrukkelijk NIET de bedoeling dat je het grootste deel van de corpus
%% van je bachelorproef in dit hoofstuk verwerkt! Dit hoofdstuk is eerder een
%% kort overzicht van je plan van aanpak.
%%
%% Maak voor elke fase (behalve het literatuuronderzoek) een NIEUW HOOFDSTUK aan
%% en geef het een gepaste titel.

%todo: beter uitschrijven:
% REQUIREMENTS-ANALYSE?
% Functionaliteiten prioriteiten MOSCOW methodologie?

\section{Marktonderzoek}
\label{sec:marktonderzoek}
%todo: vergelijking Offline -> online mogelijkheden
Eerst zal er onderzocht worden welke providers en platformen er bestaan en wat het kost om AI te gebruiken, waar vindt men de laatste informatie en ontwikkelingen van AI, hoe kunnen ze met elkaar vergeleken worden qua performantie, nauwkeurigheid en kost-efficiëntie. Ook zal er onderzocht worden welke AI en agentic tools en systemen bestaan,  welke hun capaciteiten en mogelijkheden zijn, hoe deze werken in code en uiteindelijk welke codeer-taal het best ingezet wordt om deze zelf te implementeren.

\section{Vergelijkende Studie}
\label{sec:vergelijkendestudie}
% j J

De verschillende technologieën worden grondig onderzocht om hun mogelijkheden in kaart te brengen. Op basis van de sterktes en zwaktes wordt vervolgens bepaald welke technologie het meest geschikt is als AI-agent. Mogelijks kunnen er meerdere uitwerkingen worden gevonden die elk een zinvolle taak hebben als AI agent.

\section{Opbouw Prototype AI Agent}
\label{sec:opbouwprototypeaiagent}

Met gevonden AI-technologieën zal er getracht worden een proof of concept AI agent of mischien meerdere AI agents op te zetten die kunnen ingezet worden om te helpen bij de ontwikkeling en testen van APIs, er wordt ook onderzocht voor welke taken deze best kunnen dienen om het werk van ontwikkelaars efficiënter te maken.

Ook hardware restricties, security risico's en kosten worden in acht genomen.

\section{Agent inzetten en testen}
\label{sec:agentinzettenentesten}

De gekozen technologie en gemaakte prototype(s) zal(zullen) getest worden op verschillende KPI's (Key Performance Indicators) zoals foutdetectie; productiviteit; kwaliteit; autonomie; performantie; test-kwaliteit; documentatie-kwaliteit. Hierbij kunnen verschillende LLMs ook met elkaar vergeleken worden om een optimale werking te bekomen van de uitgewerkte proof of concept AI agent.


